% !TeX root = se100_the_ontological_neutrality_theorem.tex

\section{The Impossibility of Causal and Normative Commitments in a Neutral Ontological Substrate}
\label{sec:impossibility}

Because interpretive frameworks derive conclusions
while ontologies assert commitments,
embedding framework-contestable propositions as substrate commitments
produces incompatibility with some admissible framework.
The following argument formalizes this constraint.

This section shows that the neutrality requirements
defined in Section~\ref{sec:requirements}
are incompatible with asserting causal or normative commitments
as substrate-layer ontological facts.
The argument is structural and logical rather than empirical:
it follows from the interaction
between interpretive disagreement and substrate-level ontological assertion.
All necessity and impossibility claims in this section
are relative to those requirements.

\subsection{Neutrality as Compatibility Across Admissible Frameworks}

By definition, a neutral ontological substrate must remain compatible
with multiple admissible interpretive frameworks,
even when those frameworks disagree.
Compatibility here is understood in a strong sense:
the same substrate must remain logically consistent
when extended by divergent interpretive frameworks,
without revision or retraction of substrate-layer assertions.

A framework is \emph{admissible} if it is internally consistent and
represents a legitimate interpretive stance
within the domain of application.
Admissible interpretive frameworks may disagree about
causal or normative conclusions concerning the same entities or events,
while remaining internally consistent.
For example,
a recognized legal jurisdiction, an established causal modeling tradition,
or a coherent normative theory.
Admissibility is not universality:
admissible frameworks may contradict one another,
as when two jurisdictions reach opposite legal conclusions
or two causal models attribute responsibility differently.
The set of admissible frameworks $\mathbb{F}$ is thus
characterized by internal consistency, not mutual compatibility.

This neutrality requirement immediately constrains
the Level of Abstraction (LoA)
at which the substrate,
under the neutrality definition,
must operate.
For a substrate to be neutral,
its foundational commitments must remain invariant
across all admissible interpretive frameworks.
If the substrate asserts a proposition whose truth
depends on a particular causal or normative interpretation,
then there exist admissible frameworks that derive the negation of that proposition.
Because neutrality requires the substrate to remain logically
consistent when extended by every admissible framework,
asserting such a proposition at the substrate layer
necessarily produces inconsistency with
at least one admissible framework extension.
The neutrality condition is therefore violated.

This restriction does not prohibit the representation of
causal or normative concepts themselves.
A neutral substrate may contain vocabulary for representing
causal or normative claims, provided those claims are represented as
interpretive assertions rather than asserted as substrate-layer facts.
What neutrality excludes is the elevation of a particular
causal or normative conclusion to a foundational commitment of the ontology,
since such conclusions may be rejected by other admissible frameworks.

\subsection{Normative Commitments and Interpretive Incompatibility}
\label{subsec:normative}

Normative conclusions, such as claims that
an action constituted a violation,
that an obligation applied to a given actor,
or that a rule was permissibly invoked,
are inherently framework-relative.
Their truth depends on legal regimes, institutional authorities,
temporal scope, and interpretive stance.
It is therefore a routine and expected feature
of accountability systems
that admissible frameworks disagree about normative conclusions.

Suppose an ontology asserts a normative conclusion at the substrate layer,
such as that a particular action was prohibited or
that an obligation applied to a given actor.
By doing so, the ontology privileges one normative framework over others.
Any framework that denies the asserted conclusion becomes
incompatible with the substrate,
as it cannot be layered atop it
without contradiction or revision.

This incompatibility is not a matter of
missing context or insufficient detail.
It arises from the ontological act of
asserting a contested normative conclusion as a fact.
As a result, no ontology that asserts normative conclusions
as foundational substrate facts
can, under the neutrality requirements,
satisfy interpretive non-commitment.
Moreover, when alternative normative interpretations arise,
as they inevitably do,
the substrate must be revised to restore consistency,
thereby violating extension stability.

\subsection{Causal Commitments and Model Dependence}
\label{subsec:causal}

Causal statements exhibit the same structural problem,
despite their different surface form.
Claims such as \textit{event A caused event B}
depend on background assumptions about causal mechanisms,
variable selection, counterfactual reasoning, and model scope.
Distinct causal frameworks may be equally admissible
while disagreeing about specific causal relationships.

If an ontology asserts causal relations at the substrate layer,
it necessarily commits to one causal model among many.
Frameworks that reject that model,
or that attribute causation differently,
cannot be layered onto the substrate without conflict.
As with normative commitments,
the ontology ceases to be interpretively neutral and,
under the extension stability requirement,
cannot remain stable under extension.

Importantly, this argument does not rely on
the presence of an active dispute over a specific event.
Rather, it shows that causal attribution
cannot be asserted as a substrate-layer fact
without committing the ontology
to a particular causal model.
Because causal attribution depends on background modeling assumptions,
asserting such relations at the foundational layer
constitutes a structural commitment
to a particular counterfactual or mechanistic logic.
As soon as an ontology asserts $A \to B$ as a causal fact,
it excludes any framework that treats $A$ and $B$
as merely correlated or as having a common cause $Z$.
Therefore, under the neutrality requirements,
causal conclusions cannot be asserted as substrate-layer facts:
by asserting a causal attribution,
which is a function of an interpretive model,
as a substrate-level fact,
the ontology conflates the object of observation
with its explanatory evaluation.

To be clear, \textit{pre-causal does not mean acausal}
(i.e., causation-denying).
The substrate does not deny that causation exists
or that causal reasoning is valuable.
Rather, it refrains from embedding causal conclusions
as foundational facts,
reserving causal attribution for interpretive layers
where framework assumptions are explicit
and disagreement is representable.
Causation is externalized, not eliminated.

\subsection{Reification Does Not Eliminate the Problem}
\label{subsec:reification}

One might attempt to preserve neutrality by
reifying causal or normative statements,
representing them as claims, reports, or assertions
rather than as direct ontological facts.
One might attempt to preserve neutrality by reifying causal or normative statements,
representing them as claims, reports, or assertions rather than as direct ontological facts.
Such reification is necessary for representing interpretive discourse,
but it is not sufficient to preserve neutrality.

Reification allows the ontology to represent discourse about the world
without asserting conclusions about the world itself.
This strategy aligns with recent arguments for epistemic ontologies
that represent knowledge about the world
rather than the world directly~\citep{kassel2023epistemic}.
For example, an ontology may represent that
an agent asserted that event $A$ caused event $B$
without asserting the causal relation $A \rightarrow B$
as a substrate-layer fact.
In this way, reified claims are treated as entities
available for interpretation, comparison, or evaluation,
rather than as ontological commitments.

Reification preserves interpretive non-commitment
only if the ontology refrains from simultaneously
asserting the reified content as a substrate-layer truth.
If causal or normative relations are both asserted
as facts and reified as claims, neutrality is already lost.
If they are represented solely as claims,
then they no longer function as ontological commitments
but as objects of discourse layered atop the substrate.

Thus, reification does not provide a middle ground
in which causal or normative conclusions
can be embedded without consequence.
It either externalizes interpretation,
consistent with pre-causal and pre-normative design,
or it leaves the original incompatibility intact.

The preceding analysis establishes that
causal and normative conclusions are
framework-contestable and that
neutrality requires substrate assertions
to remain compatible with every admissible framework extension;
the following theorem formalizes the resulting constraint on substrate design.

\subsection{The Ontological Neutrality Theorem}
\label{subsec:formal-proof}

The preceding observations yield the central result of this paper.

\begin{theorem}[Ontological Neutrality Theorem]
  \label{thm:neutrality}
  Let $\mathcal{S}$ be an ontology intended to function
  as a neutral substrate across diverse interpretive frameworks,
  and let $\mathbb{F}$ be the set of admissible interpretive frameworks.
  Then $\mathcal{S}$ satisfies the requirements of neutrality
  (interpretive non-commitment and extension stability)
  if and only if its foundational Level of Abstraction
  does not assert causal or normative commitments
  as substrate-level ontological facts.
\end{theorem}

Here a \emph{causal or normative commitment} refers to the assertion
of a causal relation or normative proposition as a substrate-layer fact
(e.g., that event $A$ caused event $B$, or that an action constituted a violation),
rather than the representation of such propositions as interpretive assertions or reports.

\begin{proof}
\textbf{(Only if.)}
Suppose $\mathcal{S}$ asserts a causal or normative commitment $p$
as a substrate-layer fact.
Causal and normative commitments are framework-contestable,
meaning that admissible frameworks may legitimately disagree
about whether such propositions hold.
In particular, there exist frameworks
$\mathcal{F}_1, \mathcal{F}_2 \in \mathbb{F}$
such that $\mathcal{F}_1 \vdash p$ and
$\mathcal{F}_2 \vdash \neg p$.

By the definition of neutrality (Section~\ref{sec:requirements}),
$\mathcal{S}$ must satisfy:
for all $\mathcal{F} \in \mathbb{F}$,
$\mathcal{S} \cup \mathcal{F} \not\vdash \bot$.

Since $\mathcal{S}$ asserts $p$ and
$\mathcal{F}_2 \vdash \neg p$,
it follows that
\[
\mathcal{S} \cup \mathcal{F}_2 \vdash p
\quad\text{and}\quad
\mathcal{S} \cup \mathcal{F}_2 \vdash \neg p,
\]
and therefore
\[
\mathcal{S} \cup \mathcal{F}_2 \vdash p \land \neg p.
\]
To restore consistency, either:
\begin{enumerate}
  \item[(i)] $\mathcal{F}_i$ must be excluded from $\mathbb{F}$,
        violating interpretive non-commitment; or
  \item[(ii)] $\mathcal{S}$ must be revised to remove $p$,
        violating extension stability.
\end{enumerate}
In either case, $\mathcal{S}$ fails to satisfy neutrality.

\medskip
\textbf{(If.)}
Suppose $\mathcal{S}$ refrains from asserting causal or normative commitments
at the substrate level and instead represents only
framework-invariant referential structures,
such as entities and their identity conditions.
Under this restriction,
interpretive frameworks may introduce causal or normative
conclusions as extensions without contradicting the substrate,
because those conclusions are not fixed by the foundational layer.
Thus, for all $\mathcal{F} \in \mathbb{F}$,
$\mathcal{S} \cup \mathcal{F}$ remains consistent,
satisfying both interpretive non-commitment and extension stability.

\medskip
\noindent
\textit{Scope note.}
This direction assumes that entity existence and identity conditions
are not themselves contested across admissible frameworks.
In domains where identity is interpretively contested, for example,
where frameworks disagree about whether two records denote the same entity,
or whether a particular kind of entity exists at all,
the substrate, under the neutrality definition,
cannot be neutral until identity conditions are resolved.
Such resolution is a precondition for neutrality, not a product of it.
The theorem applies to domains where identity can be fixed
prior to interpretive extension;
it does not claim that all domains admit neutral substrates.
\end{proof}

This result is not contingent on particular modeling choices
or domain assumptions.
It follows directly from the functional role
the substrate is required to play under the neutrality definition.
Neutrality across disagreement and embedded interpretive commitments
are structurally incompatible.

The next section considers the implications of this result for ontology design,
clarifying what must be excluded from foundational layers
and how interpretation, evaluation, and explanation
may be externalized without loss of accountability.
